% !TEX root = ../../Groups and Monoids in Context.tex
\section{Transformation Groups and Monoids}
We start by introducing groups and monoids in the form the most frequently appear: collections of transformations.
\begin{definition}[Transformation]
    A transformation of a set $S$ is a map $S \to S$. 
\end{definition}
\begin{remark}
    This is not a particularly standard definition of transformation, but the concept of a map from a set to itself is an important concept that deserves a name\footnote{There is standard terminology for this in category theory, but I am restricting such terminology to use for maps between groups and monoids themselves (which will come up in section \ref{sec:homomorphisms}) to avoid confusion.}.
\end{remark}
\begin{definition}[Transformation Monoid]
    A transformation monoid is a collection $M$ of transformations which contains the identity transformation and is closed under composition. That is, $\Id_S \in M$, and whenever $f,g \colon S \to S$ satisfy $f,g \in M$, we also have $f \circ g \in M$.
\end{definition}
\begin{definition}[Transformation Group]
    A transformation group $G$ is a transformation monoid such that every transformation $f \in G$ is invertible with $f^{-1} \in G$.
\end{definition}

\begin{example}[Trivial Transformation Group]
    For any set $S$, the set $\{\Id_S\}$ is a transformation monoid. And since it's only element is the identity element, which is self-inverse, it is also a transformation group. These are the trivial transformation monoids or trivial transformation groups depending on whether we care that they are groups or not.
\end{example}

\begin{example}[Caesar cipher] \textcolor{red}{\boxed{\text{POORLY WRITTEN}}}
    In example \ref{alphabet}, we defined the alphabet $\mathcal{A}$. The Caesar cipher uses a fixed key $k$ between 1 and 25 (normally) and transforms each letter by shifting it forwards in the alphabet $k$ times where each shift either moves it forwards 1 position in the alphabet or for z, sends it back to a.
    Thus, if $c_k \colon \mathcal{A} \to \mathcal{A}$ is the operation the Caesar cipher performs on the alphabet for fixed key $k$, then the set $C_{25} := \{c_k \mid k \in \llb 1, 25 \rrb \}$ is a collection of transformations on $\mathcal{A}$.

    \medskip\noindent
    It's well known (and easy to see), that for each $k \in \llb 1, 25 \rrb$, ${c_k}^{-1} = c_{26-k}$, so every element of this set is invertible, and the inverses are also in the set, but it is neither closed under composition, nor is the identity transformation an element of it, so $C_{25}$ is not a transformation group: $c_1 \circ c_{25} = \Id_{\mathcal{A}} \notin C_{25}$.

    \medskip\noindent
    If we define $c_0 := \Id_{\mathcal{A}}$ and $C_{26} := C_{25} \cup \{c_0\}$, then since this follows modular arithmetic modulo 26, we actually do get that $C_{26}$ (which contains the identity transformation by construction) is closed under composition: whenever $l,m,n \in \llb 0, 25 \rrb$ satisfy $l + m \equiv n \mod 26$, $c_l \circ c_m = c_n$.
\end{example}

\begin{example}[Full Transformation Monoids]
    Let $S$ be a set. The set of all transformations on $S$ is called the full transformation monoid on $S$ (denoted $T_S$) and is indeed a transformtion monoid since the identity transformation on $S$ is a transformation on $S$, and composing transformations on $S$ gives a transformation on $S$.
    If $S$ has at least 2 elements, it is not a transformation group. Indeed, suppose $x,y \in S$ where $x \ne y$. Then, define the transformation $T\colon S \to S$ by
    \begin{equation*}
        T(s) := \begin{cases}
            x & \text{if $s \in \{x,y\}$}, \\
            s & \text{otherwise.}
        \end{cases}
    \end{equation*}
    $T$ is an element of the full transformation monoid of $S$, but $T$ is not invertible since it is not injective: $T(x) = T(y)$ and $x \ne y$.
\end{example}

\begin{example}[Symmetric Groups]
    Let $S$ be a set. The set of all invertible transformations on $S$ is called the symmetric group on $S$, often denoted $\operatorname*{Sym}(S)$. In the cases where $S = \llb 1, n \rrb$, it is often written $S_n$ and called the symmetric group on $n$ symbols. As the names imply, these are transformation groups: The identity transformation is invertible, the composition of invertible transformations is also an invertible transformation, and the inverse of a transformation is also an invertible transformation.
\end{example}

\begin{theorem}[Generated monoids and groups]
    Let $S$ be a set and $T$ be a set of transformations on $S$. Then 
    \begin{enumerate}[label=(\alph*)]
        \item there is a minimal transformation monoid containing $T$ as a subset. It is called the transformation monoid generated by $T$ and denoted $\langle T \rangle_{\text{Mon}}$ or $\langle T \rangle$ if it is clear that we are generating a transformation monoid,
        \item If every element of $T$ is invertible, there is a minimal transformation group containing $T$ as a subset. It is called the transformation group generated by $T$ and is denoted $\langle T \rangle_{\text{Grp}}$ or $\langle T \rangle$ if it is clear that we are generating a transformation group.
    \end{enumerate}
\end{theorem}
\begin{proof}\emph{(a)}

    \smallskip\noindent
    Let $M_0 := \{\Id_S\}$, and define $M_{n+1} := \{ t \circ m \mid t \in T, m \in M_n\}$ inductively. We claim that $\langle T \rangle_{\text{Mon}} = \bigcup_{n = 0}^\infty M_n$.
    To prove this, we must show that $\bigcup_{n = 0}^\infty M_n$ is a transformation monoid and that if $M$ is another transformation monoid containing $T$, then $\bigcup_{n = 0}^\infty M_n \subseteq M$.

    \medskip\noindent
    $\Id_S \in M_0 \subseteq \bigcup_{n = 0}^\infty M_n$, so the identity transformation is an element.

    \medskip\noindent
    If $m_1,m_2 \in \bigcup_{n = 0}^\infty M_n$, then there are $n_1, n_2 \in \mathbb{N}$ with $m_1 \in M_{n_1}, m_2 \in M_{n_2}$. $m_1 = t_1\circ \cdots \circ t_{n_1}$ and $m_2 = t_1' \circ \cdots \circ t_{n_2}'$ where each $t_i, t_j' \in T$. But then $m_1 \circ m_2 = t_1 \circ \cdots \circ t_{n_1} \circ t_1' \circ \cdots \circ t_{n_2}'$, which is clearly an element of $M_{n_1 + n_2}$, and hence an element of $\bigcup_{n = 0}^\infty M_n$.

    \bigskip\noindent
    \emph{(b)}
    
    \smallskip
    Exercise \ref{generated_groups_ex}.
\end{proof}

\begin{example}[Polynomials as functions]
    Let $P$ be the set of polynomials with coefficients in $\mathbb{Z}$. The elements of $P$ can be seen as transformations of $\mathbb{Z}$. The composition of two polynomials is still a polynomial, and the identity function is a polynomial, so $P$ is a transformation monoid.
\end{example}

\begin{example}[Transformation Group of Invertible Elements]
    Let $M$ be a transformation monoid. The set $\{t \in M \mid \text{$t$ is invertible}, t^{-1} \in M \}$ is then a transformation group since we have restricted precisely to the elements which allow it to be a transformation group.
\end{example}

\begin{example}[symmetry groups]
    Let $P_n$ be a set of points on a regular $n$-gon in the plane centred at the origin, then let $D_n$ be the set of rotations and reflections of the plane which fix $P_n$ (sends any point in $P_n$ to another point in $P_n$). $D_n$ is a transformation group on the set $\mathbb{R}^2$ since each rotation and reflection is invertible, and their inverses are themselves rotations or reflections which fix $P_n$, and the identity map on $\mathbb{R}^2$ is a $0\degree$ rotation. The final thing left to check is that composing a rotation with a reflection in either order results in another rotation or reflection (since they both fix $P_n$, their composition does), and some basic linear algebra shows that is true.
\end{example}